\chapter{Expressions and Types}

\begin{multicols*}{2}
\section{Primitive types and literals}

There are four primitive types, of which three are nullary:
\verb|Integer|, \verb|Float|, \verb|Char| and the function type
\verb|->|, which is used as an infix operator of kind $* \rightarrow *
\rightarrow *$, ie. constructing a type from two other types.

Any number the written form of which contains a \verb|.| character is
assigned type \verb|Float|, any other number is assigned type
\verb|Integer|. Character literals are assigned type \verb|Char|.

\section{Tuples}

\section{Algebraic types}

Algebraic types are constructed as follows: \verb|data| $t$ \verb|=|
$c_1$ \verb|+| $c_2$ \verb|+ ... +| $c_n$. $t$ is a type name and each
$c_i$ is a type constructor which is either specified as $c$ $t_1$
$t_2 \ldots t_n$, where $c$ is a constructor name and each $t_i$ is a
type or as $c$ \verb|{| $b_1$\verb|:| $t_1$ \verb|}| \verb|{|
$b_2$\verb|:| $t_2$ \verb|} ... {| $b_n$\verb|:| $t_n$ \verb|}| where
$c$ is a constructor name, each $b_i$ is a binding name and each $t_i$
is a type. For type consistency reasons, the same binding $b$ must
have the same value across the different constructors of a type. For
instance, the following code would be illegal:
\verb|data Foo = Bar { x : Int } + Moo { x : Char }|.

An instance of an algebraic type is constructed as follows: $c$ $e_1$
$e_2 \ldots e_n$. Named elements in constructors can be accessed by
the function of the same name, for instance if we assume the following
type \verb|data Foo = Foo { x : Int }|, the following expression
\verb|x (Foo 5)| would yield \verb|5|.

Algebraic type constructors may be recursive.

\section{Functions}

\subsection{Definition}

\subsection{Application}

\section{Bindings}

\section{Conditionals}

\section{Pattern matching}

\section{Typeclasses}

\subsection{Definition}

\subsection{Instanciation}

\end{multicols*}
